LEO has experienced a vast expansion in the number of tracked space objects occupying its space in recent decades. An obvious reason for this is that space is becoming increasingly accessible. In 2026, space is being utilised as a tool for enterprise and space technology is seen in almost every sector; maps, Wi-Fi, and even healthcare all use technology that has been developed for space [6]. Despite this, a significant portion of the increased number of tracked objects in LEO comes from the detection of smaller objects that were previously undetected and therefore untracked. These objects are typically classified as orbital debris, which can include fragments of spacecraft, and inactive satellites, and accounts for 83\% of all tracked space objects [7]. Space exploration has not peaked yet either. It is predicted that 20,000 more satellites will be launched into LEO in the next decade [3]. 
Due to this growth, our orbits have become heavily congested and are constantly getting busier, causing satellite owners/operators (O/Os) to pour millions of pounds into COLA to protect increasingly valuable orbital infrastructure from a mission-ending collision. In the SSA industry, a number of processes are currently in place to protect the satellites in orbit from potential conjunctions. As a result, considerable research has focussed on improving the safety. Manfletti et al. describes the current situation in terms of its financial cost to O/Os—up to \$15M in a single year to protect a 300-satellite constellation (which would be on the smaller end of constellations) [8]. However, even this cost is conservative. 
For continuation…
